\section{Computation Graphs, Chain Rule, and Backpropagation}

Neural networks are built by composing many simple operations. Once such a composition is written down, training requires one central computational task:

\begin{quote}
Given a loss function, how do we compute its derivatives with respect to all parameters efficiently?
\end{quote}

The answer is \emph{backpropagation}, which is best understood not as a mysterious separate algorithm, but as a systematic application of the chain rule to a computation graph. This section gives the formal language for that viewpoint and works through a full two-layer example in detail.

\subsection*{Computation graphs}

A neural network is naturally described as a directed acyclic graph of intermediate quantities. Each node stores a value produced from earlier nodes, and the final node stores the loss.

\begin{definition}[Computation graph]
A \emph{computation graph} is a directed acyclic graph whose nodes represent intermediate variables and whose edges indicate functional dependence. Each non-input node is obtained by applying an elementary operation to its parent nodes. A scalar loss $L$ is designated as an output node, and differentiation is performed with respect to this final scalar quantity.
\end{definition}

In neural-network training, the forward pass computes all node values from input to loss, while the backward pass computes derivatives of the loss with respect to intermediate quantities and parameters.

\begin{figure}[tbp]
\centering
\begin{tikzpicture}[
    box/.style={draw, rounded corners, fill=gray!10, inner sep=6pt, minimum width=1.9cm, minimum height=0.9cm},
    pbox/.style={draw, rounded corners, fill=gray!05, inner sep=4pt, minimum width=1.4cm, minimum height=0.8cm},
    arrow/.style={-{Latex[length=2.2mm]}, thick},
    every node/.style={align=center}
]
\node[box] (x) at (0,0) {$x$};
\node[box] (z1) at (2.7,0) {$z^{(1)}=W_1x+b_1$};
\node[box] (h) at (5.2,0) {$h=\sigma(z^{(1)})$};
\node[box] (z2) at (7.8,0) {$z^{(2)}=W_2h+b_2$};
\node[box] (yhat) at (10.3,0) {$\hat y$};
\node[box] (L) at (12.6,0) {$L=\ell(\hat y,y)$};

\node[pbox] (W1) at (2.7,1.6) {$W_1,b_1$};
\node[pbox] (W2) at (7.8,1.6) {$W_2,b_2$};
\node[pbox] (y) at (12.6,1.6) {$y$};

\draw[arrow] (0.95,0) -- (1.55,0);
\draw[arrow] (3.85,0) -- (4.25,0);
\draw[arrow] (6.15,0) -- (6.85,0);
\draw[arrow] (8.75,0) -- (9.35,0);
\draw[arrow] (11.15,0) -- (11.75,0);

\draw[arrow] (W1) -- (z1);
\draw[arrow] (W2) -- (z2);
\draw[arrow] (y) -- (L);

\end{tikzpicture}
\caption{A computation graph for a two-layer network. The forward pass computes intermediate variables from left to right; backpropagation traverses the same graph in reverse to compute derivatives of the loss with respect to all parameters.}
\end{figure}

\subsection*{The chain rule as the engine of gradient computation}

The basic principle is the chain rule. In one variable, if
\[
u=g(v),\qquad v=h(t),
\]
then
\[
\frac{du}{dt}=\frac{du}{dv}\frac{dv}{dt}.
\]

In multivariable settings, the same idea propagates derivatives through compositions of vector-valued maps. If
\[
L=L(z),\qquad z=z(\theta),
\]
then
\[
\nabla_\theta L
=
\left(\frac{\partial z}{\partial \theta}\right)^\top \nabla_z L.
\]

Backpropagation is the organized repeated use of this principle across an entire computation graph.

\subsection*{Forward values and backward adjoints}

For each node $v$ in the computation graph, one may associate an \emph{adjoint}
\[
\bar v := \frac{\partial L}{\partial v},
\]
that is, the derivative of the final loss with respect to that node.

The forward pass computes the values
\[
v_1,v_2,\dots,v_N.
\]
The backward pass computes the corresponding adjoints
\[
\bar v_1,\bar v_2,\dots,\bar v_N
\]
starting from
\[
\bar L = \frac{\partial L}{\partial L}=1
\]
and moving backward through the graph.

This gives the conceptual structure of reverse-mode differentiation:

\begin{itemize}
    \item compute all intermediate values once;
    \item start from the loss;
    \item propagate sensitivities backward using local derivatives.
\end{itemize}

\subsection*{Backpropagation as reverse-mode differentiation}

Suppose a node $u$ feeds into several later nodes. Then changing $u$ affects the loss through all downstream paths. The total derivative therefore sums the contributions from all children:
\[
\frac{\partial L}{\partial u}
=
\sum_{v:\,u\to v}
\frac{\partial L}{\partial v}\frac{\partial v}{\partial u}.
\]
This is the local recursion underlying backpropagation.

\begin{proposition}[Reverse-mode differentiation is efficient]
For a scalar loss defined on a computation graph, reverse-mode differentiation computes the gradient of the loss with respect to all parameters at a computational cost comparable to one forward pass, up to constant factors.
\end{proposition}

\begin{proof}[Proof idea]
A naive approach would differentiate the loss separately with respect to each parameter, repeatedly recomputing many shared intermediate derivatives. Reverse-mode differentiation avoids this redundancy.

During the forward pass, each node value is computed once and stored. During the backward pass, each edge of the graph is traversed in reverse order, and each local derivative contributes once to the adjoint of its parent node. Because shared downstream computations are reused rather than recomputed separately for each parameter, the total cost scales with the size of the graph rather than with the number of parameters times the graph size.

Under the usual assumption that each elementary operation and its local derivative can be evaluated at comparable cost, the total backward cost is bounded by a constant multiple of the forward cost. Thus all parameter gradients are obtained simultaneously at cost comparable to one forward evaluation, up to constant factors.
\end{proof}

This proposition explains why deep learning is computationally feasible at all. A network may contain millions or billions of parameters, yet backpropagation does not compute each derivative from scratch.

\subsection*{A formal worked example: a two-layer network}

We now work through a full example. Let
\[
x\in\mathbb R^d,
\qquad
z^{(1)}=W_1x+b_1 \in \mathbb R^m,
\qquad
h=\sigma(z^{(1)}) \in \mathbb R^m,
\]
\[
z^{(2)}=W_2h+b_2 \in \mathbb R^k,
\qquad
\hat y=z^{(2)}.
\]
For simplicity, take squared loss
\[
L=\frac12\|\hat y-y\|^2.
\]

\begin{example}[Backpropagation through a two-layer network]
For the network above, the parameter gradients are:
\[
\frac{\partial L}{\partial \hat y}=\hat y-y,
\]
\[
\frac{\partial L}{\partial W_2}
=
(\hat y-y)\,h^\top,
\qquad
\frac{\partial L}{\partial b_2}
=
\hat y-y,
\]
\[
\delta^{(1)}
:=
\frac{\partial L}{\partial z^{(1)}}
=
\left(W_2^\top(\hat y-y)\right)\odot \sigma'(z^{(1)}),
\]
\[
\frac{\partial L}{\partial W_1}
=
\delta^{(1)}x^\top,
\qquad
\frac{\partial L}{\partial b_1}
=
\delta^{(1)}.
\]
\end{example}

\begin{proof}
Since
\[
L=\frac12\|\hat y-y\|^2
=
\frac12\sum_{j=1}^k (\hat y_j-y_j)^2,
\]
we have
\[
\frac{\partial L}{\partial \hat y}=\hat y-y.
\]
Because
\[
\hat y=z^{(2)}=W_2h+b_2,
\]
the derivative with respect to the affine parameters of the second layer is
\[
\frac{\partial L}{\partial W_2}
=
\frac{\partial L}{\partial z^{(2)}}\frac{\partial z^{(2)}}{\partial W_2}
=
(\hat y-y)h^\top,
\]
and similarly
\[
\frac{\partial L}{\partial b_2}=\hat y-y.
\]

Next, we propagate the derivative backward to the hidden layer:
\[
\frac{\partial L}{\partial h}
=
W_2^\top(\hat y-y).
\]
Since
\[
h=\sigma(z^{(1)}),
\]
applied componentwise, the chain rule gives
\[
\delta^{(1)}
=
\frac{\partial L}{\partial z^{(1)}}
=
\frac{\partial L}{\partial h}\odot \sigma'(z^{(1)})
=
\left(W_2^\top(\hat y-y)\right)\odot \sigma'(z^{(1)}).
\]
Finally, because
\[
z^{(1)}=W_1x+b_1,
\]
the gradients of the first-layer parameters are
\[
\frac{\partial L}{\partial W_1}
=
\delta^{(1)}x^\top,
\qquad
\frac{\partial L}{\partial b_1}
=
\delta^{(1)}.
\]
This is exactly the backpropagation formula for the two-layer network.
\end{proof}

The pattern here is the one that will recur throughout the book:

\begin{enumerate}
    \item compute forward activations;
    \item compute loss;
    \item propagate error signals backward;
    \item use those signals to form parameter gradients.
\end{enumerate}

\subsection*{Forward and backward signal flow}

The forward pass carries \emph{representations} from input to output. The backward pass carries \emph{sensitivity information} from the loss back to earlier layers.

\begin{figure}[tbp]
\centering
\begin{tikzpicture}[
    box/.style={draw, rounded corners, fill=gray!10, inner sep=6pt, minimum width=1.8cm, minimum height=0.8cm},
    farrow/.style={-{Latex[length=2.2mm]}, thick},
    barrow/.style={-{Latex[length=2.2mm]}, thick, dashed},
    every node/.style={align=center}
]
\node[box] (x) at (0,0) {$x$};
\node[box] (h1) at (2.8,0) {$z^{(1)},h$};
\node[box] (h2) at (5.6,0) {$z^{(2)},\hat y$};
\node[box] (L) at (8.4,0) {$L$};

\draw[farrow] (0.95,0) -- (1.85,0);
\draw[farrow] (3.75,0) -- (4.65,0);
\draw[farrow] (6.55,0) -- (7.45,0);

\draw[barrow] (7.45,-0.55) -- (6.55,-0.55);
\draw[barrow] (4.65,-0.55) -- (3.75,-0.55);
\draw[barrow] (1.85,-0.55) -- (0.95,-0.55);

\node at (4.2,0.7) {\small forward: activations / representations};
\node at (4.2,-1.1) {\small backward: gradients / error signals};
\end{tikzpicture}
\caption{Forward and backward signal flow. The network computes activations from left to right, then propagates gradient information from the loss back toward earlier layers.}
\end{figure}

This dual flow explains a deep conceptual point: training requires both \emph{evaluation} and \emph{credit assignment}. The forward pass evaluates the model on an input; the backward pass assigns responsibility for the loss to earlier computations and parameters.

\subsection*{Layerwise form of backpropagation}

For a multilayer feedforward network with layer equations
\[
z^{(\ell)}=W_\ell h^{(\ell-1)}+b_\ell,
\qquad
h^{(\ell)}=\sigma_\ell(z^{(\ell)}),
\]
the backpropagation recursion has the generic form
\[
\delta^{(\ell)}
=
\left(W_{\ell+1}^\top \delta^{(\ell+1)}\right)\odot \sigma_\ell'(z^{(\ell)}),
\]
together with parameter gradients
\[
\frac{\partial L}{\partial W_\ell}
=
\delta^{(\ell)}\bigl(h^{(\ell-1)}\bigr)^\top,
\qquad
\frac{\partial L}{\partial b_\ell}
=
\delta^{(\ell)}.
\]

This formula compresses much of the practical content of backpropagation. Each layer receives an incoming error signal from above, transforms it through the transpose of the weight matrix and the derivative of the activation, and then uses the result to compute its own parameter gradients.

\subsection*{Backpropagation versus learning}

Backpropagation computes gradients. It does not itself specify how parameters are updated after those gradients are known.

Once gradients have been computed, one still needs an optimization rule, such as gradient descent or stochastic gradient descent:
\[
W_\ell \leftarrow W_\ell - \eta \frac{\partial L}{\partial W_\ell},
\qquad
b_\ell \leftarrow b_\ell - \eta \frac{\partial L}{\partial b_\ell}.
\]
So backpropagation belongs to the differentiation side of learning, not the optimization side.

\commonconfusion{Backpropagation is not the learning rule. It is the gradient-computation procedure. The learning rule is the parameter-update method applied after gradients have been computed, such as gradient descent, SGD, or Adam. Confusing these two levels hides an important conceptual distinction: one procedure computes sensitivities, the other decides how to move in parameter space.}

\whattoremember{A computation graph is a directed acyclic graph of intermediate variables ending in a scalar loss. Backpropagation is reverse-mode differentiation on this graph: it applies the chain rule efficiently by propagating adjoints backward from the loss. Its key computational advantage is that it computes all parameter gradients at cost comparable to one forward pass up to constant factors. In practice, training a neural network means combining forward evaluation, backward gradient computation, and then a separate optimization update.}

\subsection*{Exercises}

\begin{exercise}
Consider the scalar computation
\[
u=ax+b,\qquad v=\sigma(u),\qquad L=\frac12(v-y)^2.
\]
Compute
\[
\frac{\partial L}{\partial a},
\qquad
\frac{\partial L}{\partial b}
\]
using the chain rule.
\end{exercise}

\begin{exercise}
For the two-layer network
\[
z^{(1)}=W_1x+b_1,\qquad h=\sigma(z^{(1)}),\qquad \hat y=W_2h+b_2,
\]
derive
\[
\frac{\partial L}{\partial W_2}
\]
under squared loss
\[
L=\frac12\|\hat y-y\|^2.
\]
\end{exercise}

\begin{exercise}
Using the same two-layer network, derive the hidden-layer error signal
\[
\delta^{(1)}=\frac{\partial L}{\partial z^{(1)}}.
\]
Explain where the Hadamard product with $\sigma'(z^{(1)})$ comes from.
\end{exercise}

\begin{exercise}
Why is reverse-mode differentiation especially well suited to neural-network training, where the output of interest is a scalar loss but the number of parameters is very large?
\end{exercise}

\begin{exercise}
Explain in words the difference between the forward pass and the backward pass in a neural network.
\end{exercise}

\begin{exercise}
A student says, ``Backpropagation trains the network.'' What is incomplete or inaccurate about this statement?
\end{exercise}

\subsection*{Notes and references}

This section formalizes one of the central computational ideas of deep learning: gradients in layered models are obtained efficiently by reverse-mode differentiation on a computation graph. The conceptual ingredients are simple---composition, the chain rule, stored intermediate values, and recursive sensitivity propagation---but together they make the training of large neural networks computationally practical. Later sections will build on this by studying when gradient-based training works well, when it becomes unstable, and how architecture, initialization, normalization, and residual structure influence signal propagation through depth.